\documentclass[journal]{IEEEtran}
\usepackage[T1]{fontenc}
\usepackage{url}
\title{Your Manuscript Title}
\author{Your Name}
\begin{document}
\maketitle
\begin{abstract}
Replace this placeholder with the abstract of your real manuscript.
\end{abstract}
\begin{IEEEkeywords}
Replace these keywords with your manuscript's keywords.
\end{IEEEkeywords}
\section{Getting Started}
This is a minimal IEEE starter, not a manuscript ready for submission.
Replace \texttt{original/paper.tex} with your real manuscript and
\texttt{original/paper.bib} with your real bibliography. Keep the filenames
\texttt{paper.tex} and \texttt{paper.bib} so the supplied build rules work.
Replace the title, author, abstract, keywords, and all placeholder sections.
The citation below is fictional; remove it when you add your real references~\cite{placeholder}.
\section{Introduction}
Replace this section with your manuscript's introduction and add the sections you need.
Keep included source files and figures inside \texttt{original/}.
Run \texttt{make} from the project root to build \texttt{original/paper.pdf}.
\section{Preparing Revisions}
After submitting the original manuscript, freeze its sources and run
\texttt{make new-round}. This creates \texttt{revised1} and \texttt{reply1}.
Edit the revised manuscript and fill in the rebuttal metadata, comments, and answers.
Run \texttt{make} to build the documents or \texttt{make bundle} to append comparisons.
Use \texttt{make new-round} again for each later review round.
\section{Conclusion}
Replace this section with your manuscript's conclusions.
\bibliographystyle{IEEEtran}
\bibliography{paper}
\end{document}
